\subsection{An unconditional three-class bound for the degree-fifty-six field}

We now remove the finite class-group hypothesis by exhibiting enough
principal prime ideals to exclude an unramified cyclic cubic extension.
The method is the explicit-formula class-number argument: principal
primes split in an unramified abelian extension, and their contributions
raise the lower bound for its discriminant.

\begin{proposition}\label{prop:septic-three-class}
For the degree-$56$ field $M$ of
Proposition~\ref{prop:lm-septic-degree56},
\[
 \operatorname{Cl}(M)[3]=0.
\]
\end{proposition}
\begin{proof}
Suppose otherwise. Class field theory gives an unramified cyclic
cubic extension $T/M$. Since the extension has odd degree, its real
places remain real. Hence
\[
 [T:\mathbb Q]=168,\qquad (r_1(T),r_2(T))=(6,81),\qquad
 \operatorname{rd}(T)=\operatorname{rd}(M)
   =3^{55/56}7^{62/56}<25.366.
 \tag{$*$}\label{eq:septic-hilbert-rd}
\]

Here is the finite algebraic input to the discriminant calculation.
The field can be constructed from $\alpha^7=3$ and
$w=3z$, whose monic polynomial is
\[
 w^8-441w^6+18522w^4-28588707=0.
\]
The exact maximal-order computation verifies its degree, signature
and discriminant. For $592$ distinct prime ideals of $M$ of norm
at most $5000$, the accompanying certificate gives an algebraic
integer generating each ideal. It verifies the equality of the
two ideal lattices in Hermite normal form and the exact norm of
the generator. There are $204$ distinct norms, with largest norm
$4973$. The replay verifies irreducibility of the degree-eight
and degree-$56$ polynomials, the two subfield embeddings, the
maximal order, the residue degrees and primality of each ideal,
and absence of repeated ideals. In particular the input includes one ideal of norm $3$,
two of norm $5$, and three of norm $25$.
The verification uses only arithmetic in the certified maximal
order and the explicit generators; it does not assume a computed
class group or the generalized Riemann hypothesis.
Write $c_q$ for the number of these verified ideals of norm $q$.
Each splits completely in $T$, so $T$ has at least $3c_q$ primes
of norm $q$.

For completeness we specify the analytic calculation, including
the error bounds. Put $g(t)=1+\cos(\pi t)$ for $|t|\le1$ and
$g(t)=0$ otherwise. Its normalized autocorrelation is the even
function $H$, supported on $[-2,2]$, whose expression for
$0\le t\le2$ is
\[
 H(t)=\frac{(2-t)(1+\tfrac12\cos\pi t)}3
       +\frac{\sin\pi t}{2\pi}.
\]
It has $H(0)=1$, $0\le H\le1$, and nonnegative Fourier
transform. Set $b=14$ and
\[
 F(x)=\frac{H(x/b)}{\cosh(x/2)}.
\]
This is an unconditional admissible kernel in the explicit formula \citep[Eqs.~(2.3)--(2.4), p.~122]{Odlyzko1990}.
Define
\[
 I=\int_0^\infty\frac{1-F(x)}{2\sinh(x/2)}\,dx,
 \qquad
 A=\int_0^{2b}\frac{H(x/b)}{2\cosh^2(x/2)}\,dx.
\]
The real-place term is $A$, since
$\int_0^\infty(2\cosh(x/2))^{-1}\,dx=\pi/2$.
Also $\int_0^2H(t)\,dt=2/3$, so the pole term is
$4\int_0^\infty F(x)\cosh(x/2)\,dx=8b/3$.
Discarding the nonnegative zero contributions and all unused
prime ideals in the explicit formula therefore gives
\[
\begin{split}
 \log\operatorname{rd}(T)\ge{}&
 \log(8\pi)+\gamma-I+\frac A{28}-\frac29\\
 &+\frac1{56}\sum_q c_q
   \sum_{\substack{m\ge1\\m\log q\le28}}
       \frac{4\log q\,H(m\log q/14)}{q^m+1}.
\end{split}
\tag{$**$}\label{eq:septic-explicit-lower}
\]
The factor $1/56$ is $3/168$: every verified prime of $M$
gives three primes in $T$. This calculation does not require
$M/\mathbb Q$ or $T/\mathbb Q$ to be Galois.

The finite sum in~\eqref{eq:septic-explicit-lower}, before division by $56$, has the certified
value
\[
 19.51514734319693587857082907544961
 \quad\text{with error less than }10^{-31}.
\]
The integrals are enclosed using $192$-bit ball arithmetic.
For $\varepsilon=1/16384$, the integral defining $I$ from
$\varepsilon$ to $28$ is evaluated by rigorous complex-ball
quadrature. Its tail is exactly $-\log\tanh7$.
The omitted integral from $0$ to $\varepsilon$ lies between
$0$ and $\varepsilon^2/2$: indeed $H'(0)=0$, $|H''|\le6$,
and $1-\operatorname{sech}(x/2)\le x^2/8$ imply that its
integrand is at most $x$. Using this upper endpoint bound for
$I$ gives
\[
 I<0.726948814142097,
 \qquad A>0.973804608359249.
\]
Substitution in~\eqref{eq:septic-explicit-lower}, with outward rounding throughout, yields
\[
 \operatorname{rd}(T)>25.4185>25.366,
\]
contradicting~\eqref{eq:septic-hilbert-rd}. Thus no such cubic extension exists, and
$\operatorname{Cl}(M)[3]=0$.
\end{proof}

A second evaluation with $b=27/2$ and $256$-bit arithmetic gives
the compatible lower bound $25.4115>25.366$, using the same exact
principal ideals. The proof establishes the required three-class
vanishing and does not need the exact class number of $M$.

By Proposition~\ref{prop:lm-septic-degree56}, the residual field
$K$ also has $\operatorname{Cl}(K)[3]=0$. Applying
Corollary~\ref{cor:lm-septic-class-gate} completes the septic
case, and hence all three families, of
Theorem~\ref{thm:unrestrictedthreefamiliesintro}.
